

\section{Rain gauge dataset} \label{app:rainfall}

Describe the first dataset here: where it came from, how it was
processed, and any caveats a reader would need in order to reuse
it.

\lipsum[21]

% [ht] instead of the default [t]: an appendix float left to
% drift to the top of the page can end up printed above its own
% section heading. This is the optional position argument in use.
\begin{TableM}[ht]{%
  Structure of the rain gauge dataset: one record per gauge and
  day.%
}{tbl:app:rainfall}
\scriptsize
\sffamily
\begin{tabularx}{\linewidth}{LLL}
\toprule
\textbf{Field} & \textbf{Unit} & \textbf{Description} \\
\midrule
\texttt{gauge\_id} & ---    & Station identifier \\
\texttt{date}      & ---    & Observation date, ISO 8601 \\
\texttt{rainfall}  & mm     & Daily accumulated rainfall \\
\texttt{quality}   & ---    & Quality flag, 0--3 \\
\bottomrule
\end{tabularx}
\end{TableM}

\lipsum[22]

\section{Streamflow dataset} \label{app:streamflow}

Describe the second dataset here, following the same pattern.
\Cref{fig:app:forecast} shows that figures and tables inside an
appendix pick up lettered numbering automatically.

\lipsum[23]

\FigureM[ht]{figure_66}{%
  Supplementary precipitation forecasts for the April--May 2024
  event. See \Cref{fig:forecast} for model and date details.%
}{fig:app:forecast}

\lipsum[24]
